\section*{Цель работы}
Практическое ознакомление со схемами включения операционного усилителя (ОУ), исследование его основных характеристик и параметров: передаточной характеристики, амплитудно-частотной характеристики (АЧХ), напряжения смещения и скорости нарастания выходного напряжения.

\section*{Задачи исследования}
\begin{enumerate}
    \item Измерить напряжение смещения $U_{\text{см}}$ на выходе ОУ.
    \item Исследовать передаточную характеристику ОУ $U_{\text{вых}} = f(U_{\text{вх}})$ при различных сопротивлениях нагрузки.
    \item Снять АЧХ усилителя, определить частоту единичного усиления $f_{\text{ЕУ}}$ и верхнюю границу полосы усиления максимальной мощности $f_{\text{УММ}}$.
    \item Определить максимальную скорость нарастания выходного напряжения $v$.
\end{enumerate}

\section*{Теоретическое введение}
Операционный усилитель (ОУ) --- это многокаскадный усилитель постоянного тока, предназначенный для выполнения различных операций над аналоговыми сигналами. ОУ имеет два входа: \textit{инвертирующий} (обозначается знаком «-») и \textit{неинвертирующий} (обозначается знаком «+»), и один выход.

\begin{figure}[H]
    \centering
    \includegraphics[width=0.4\textwidth]{figs/opamp_symbol.pdf}
    \caption{Условное графическое обозначение ОУ}
    \label{fig:opamp_symbol}
\end{figure}

Ключевыми характеристиками ОУ являются:
\begin{itemize}
    \item \textit{Передаточная характеристика} --- зависимость выходного напряжения $U_{\text{вых}}$ от входного $U_{\text{вх}}$ при частоте $f=0$. В линейном режиме работы эта зависимость близка к прямой.
    \item \textit{Амплитудно-частотная характеристика (АЧХ)} --- зависимость модуля коэффициента усиления $|K|$ от частоты $f$. Для ОУ характерно снижение усиления на высоких частотах.
    \item \textit{Входное и выходное сопротивление} ($R_{\text{вх}}$, $R_{\text{вых}}$). Идеальный ОУ имеет $R_{\text{вх}} \to \infty$ и $R_{\text{вых}} \to 0$.
\end{itemize}

Экспериментальная установка собрана на базе лабораторного макета, схема которого представлена на \cref{fig:main_circuit}.

\begin{figure}[H]
    \centering
    \includegraphics[width=0.9\textwidth]{figs/main_circuit.pdf}
    \caption{Принципиальная схема лабораторного макета}
    \label{fig:main_circuit}
\end{figure}

\section*{Порядок выполнения работы}

\subsection*{1. Измерение напряжения смещения ($U_{\text{см}}$)}
\begin{itemize}
    \item \textit{Задача:} Определить напряжение на выходе ОУ при отсутствии входного сигнала, обусловленное внутренними асимметриями схемы.
    \item \textit{Порядок действий:}
        \begin{enumerate}
            \item Собрать схему инвертирующего усилителя. Подключить резистор $R_1$ к инвертирующему входу ОУ (клеммы 2--9). Неинвертирующий вход (клемма 18) соединить с землей. В качестве цепи обратной связи использовать резистор $R_{12}$ (соединить клеммы 9--21 и 28--21).
            \item К выходу ОУ (клемма 28) подключить вольтметр постоянного тока.
            \item Измерить выходное напряжение $U_{\text{вых}}$. Результат занести в \cref{tab:offset_measurement}.
        \end{enumerate}
    \item \textit{Примечание:} Расчет напряжения смещения, приведенного ко входу, производится по формуле $U_{\text{см}} = -U_{\text{вых}} \cdot (R_1/R_{12})$ и выполняется на этапе обработки данных.
\end{itemize}

\subsection*{2. Исследование передаточной характеристики}
\begin{itemize}
    \item \textit{Задача:} Снять зависимость $U_{\text{вых}} = f(U_{\text{вх}})$ для двух различных нагрузок.
    \item \textit{Порядок действий:}
        \begin{enumerate}
            \item Собрать схему: подключить к инвертирующему входу (клемма 9) регулируемый источник постоянного напряжения через резистор $R_2$ (клеммы 3--9). В качестве обратной связи использовать $R_{12}$ (клеммы 9--21, 28--21). К выходу ОУ (клемма 28) подключить нагрузочный резистор $R_{14}$ (клеммы 28--24) и вольтметр постоянного тока.
            \item Установить на входе напряжение $U_{\text{вх}} = 0~\text{В}$. Вращая потенциометр $R_{18}$, добиться нулевого напряжения на выходе ОУ ($U_{\text{вых}} = 0~\text{В}$).
            \item Изменяя входное напряжение $U_{\text{вх}}$ в диапазоне от $-2~\text{В}$ до $+2~\text{В}$ с шагом $0.5~\text{В}$, фиксировать соответствующие значения $U_{\text{вых}}$. Результаты занести в \cref{tab:transfer_characteristic}.
            \item Заменить нагрузку $R_{14}$ на $R_{15}$ (клеммы 28--25) и повторить измерения из п. 3.
        \end{enumerate}
\end{itemize}

\subsection*{3. Исследование АЧХ}
\begin{itemize}
    \item \textit{Задача:} Измерить АЧХ и определить частоту единичного усиления $f_{\text{ЕУ}}$ и частоту усиления максимальной мощности $f_{\text{УММ}}$.
    \item \textit{Порядок действий:}
        \begin{enumerate}
            \item Собрать схему: к инвертирующему входу ОУ через резистор $R_2$ подключить генератор гармонических сигналов. Обратная связь --- резистор $R_{12}$. К выходу (клемма 28) подключить осциллограф и вольтметр переменного тока.
            \item Установить на выходе генератора напряжение $U_{\text{вх}} = 20~\text{мВ}$. Изменяя частоту от $1~\text{кГц}$ в сторону увеличения, измерять $U_{\text{вых}}$. Заполнить \cref{tab:afc_data}.
            \item Определить частоту $f_{\text{ЕУ}}$, при которой выходное напряжение становится равным входному ($U_{\text{вых}} = U_{\text{вх}} = 20~\text{мВ}$). Занести значение в \cref{tab:characteristic_freqs}.
            \item Заменить резистор обратной связи $R_{12}$ на $R_{13}$ (клеммы 9--22, 28--22). Установить на генераторе частоту $1~\text{кГц}$ и увеличить входное напряжение до получения на выходе неискаженного сигнала амплитудой $U_{\text{вых}} = 8~\text{В}$.
            \item Увеличивая частоту, зафиксировать значение $f_{\text{УММ}}$, при котором на осциллографе появляются заметные искажения формы выходного сигнала. Занести значение в \cref{tab:characteristic_freqs}.
        \end{enumerate}
\end{itemize}

\subsection*{4. Определение скорости нарастания}
\begin{itemize}
    \item \textit{Задача:} Измерить максимальную скорость изменения выходного напряжения.
    \item \textit{Порядок действий:}
        \begin{enumerate}
            \item Использовать схему из предыдущего пункта (ОС через $R_{13}$).
            \item Переключить генератор в режим формирования прямоугольных импульсов (меандр) с частотой $100~\text{кГц}$ и амплитудой $10~\text{В}$.
            \item С помощью осциллографа измерить длительность фронта выходного импульса $\tau_{\text{фр}}$. Результат занести в \cref{tab:slew_rate}.
        \end{enumerate}
    \item \textit{Примечание:} Расчет скорости нарастания производится по формуле $v \approx \Delta U / \tau_{\text{фр}}$, где $\Delta U$ — полная амплитуда выходного сигнала (пик-пик).
\end{itemize}

% ----------
% Соответствие изображений из методического пособия и label в документе
% Рис. 6.1 -> fig:opamp_symbol
% Рис. 6.4 -> fig:main_circuit
% ----------